\section{Topology-Synchronous Decoupled Observation Framework}
\label{sec:observations}

\subsection{Natural Bidirectional Excitation}

At a valid $u:0\rightarrow1$ transition, the capacitor current changes from $i_C^-=i_1^-$ to $i_C^+=-i_2^+$. Hence
\begin{equation}
\Delta i_C=i_C^+-i_C^-=-(i_1^-+i_2^+).
\label{eq:current_commutation}
\end{equation}
Because the internal ideal-capacitor voltage cannot jump, $v_C^+=v_C^-$. The terminal voltages immediately before and after the edge satisfy
\begin{align}
v_T^-&=v_C+r_Ci_1^-,\\
v_T^+&=v_C-r_Ci_2^+,
\end{align}
which gives the ideal ESR signature
\begin{equation}
v_T^- - v_T^+=r_C(i_1^-+i_2^+).
\label{eq:ideal_edge}
\end{equation}
Thus, the topology transition provides a large, signed current change while cancelling the continuous internal voltage.

\subsection{Timestamp-Reconstructed ESR Observation}

Direct edge samples are contaminated by acquisition skew, finite bandwidth, and switching transients. The proposed measurement therefore fits local linear models to timestamped points in safe pre- and post-edge windows and extrapolates both models to a common switching timestamp. With reconstructed intercepts $\hat v_T^-$ and $\hat v_T^+$, the accepted ESR row is
\begin{equation}
z_R=\hat v_T^- -\hat v_T^+
=k_R I_\Sigma r_C+\nu_R,
\quad I_\Sigma=i_1^-+i_2^+,
\label{eq:esr_observation}
\end{equation}
where $k_R$ is a training-locked calibration coefficient and $\nu_R$ collects voltage, timestamp, fitting, and residual front-end errors. The reconstruction estimates the discontinuity at one physical instant instead of subtracting samples taken at different offsets.

Simulation calibration and blind health identification use disjoint cases. Four calibration profiles---nominal, high duty/load, high noise, and nominal parasitic Model B---form the accepted-edge ratio
\begin{equation}
\kappa_m=\frac{z_{R,m}}{I_{\Sigma,m}r_{C,\mathrm{ref},m}},
\qquad
\hat k_R=\operatorname{median}_m\kappa_m,
\label{eq:kr_calibration}
\end{equation}
where $r_{C,\mathrm{ref},m}$ is the planted simulation reference ESR. The resulting training-locked coefficient is $k_R=0.9772$. The manuscript reports $k_R$ to four decimal places; the full-precision value used in all regenerated results is retained in Supplementary Note S2 and the source-traceability record.

The $\pm0.006$ spread represents case-to-case variation in the calibrated measurement-chain gain. For each simulated case, the realized coefficient is assumed known after calibration and is supplied identically to all compared estimators through $h_R=k_RI_\Sigma$. It therefore tests calibrated unit-to-unit variability rather than residual unknown calibration error; $k_R$ is never updated from blind residuals. A hardware realization would require commissioning with an independently characterized reference capacitor under the final measurement chain.

Thus, $k_R$ is known and fixed within each identification interval; otherwise the edge row identifies only the product $k_Rr_C$. The smaller evaluated-envelope value $k_{R,\min}^{\mathrm{eval}}=0.957654$ is used only for a conservative information bound and does not resolve this scale ambiguity.

\subsection{Safe-Window Charge Observation}

For an interval $[t_a,t_b]$ wholly inside one switching state, with the charge-safe-window length denoted by $T_{w,C}=t_b-t_a$, define
\begin{equation}
q=\int_{t_a}^{t_b}i_C(t)\,dt.
\label{eq:charge}
\end{equation}
Integration of~\eqref{eq:cap_port} yields
\begin{equation}
\Delta v_T=\frac{q}{C}+r_C\Delta i_C.
\label{eq:integral_port}
\end{equation}
Using the most recent accepted ESR estimate gives the capacitance pseudo-measurement
\begin{equation}
z_C=\Delta v_T-\hat r_C\Delta i_C
=\frac{q}{C_b}\alphab+\nu_C.
\label{eq:c_observation}
\end{equation}
The charge row is evaluated away from the switching edge and is capacitance-dominant after ESR compensation. ESR-estimation uncertainty and residual correlation enter the conditional measurement covariance.

More explicitly, first-order substitution of measured endpoint current, voltage, charge, and the preceding ESR estimate gives
\begin{multline}
z_C=\frac{q}{C_b}\alphab+(r_C-\hat r_C)\Delta i_C
+\nu_{\Delta v}-\hat r_C\nu_{\Delta i}+\nu_q,
\label{eq:c_mismatch}
\end{multline}
where $\nu_q$ is the voltage-equivalent charge-integration error. In the evaluated implementation its unscaled charge variance is denoted $R_q$, so that $\operatorname{var}(\nu_q)=\alphab^2R_q$. Treating the preceding ESR error as a zero-mean conditional uncertainty with variance $P_{rr}$ yields
\begin{equation}
R_C\simeq R_{\Delta v}+(\Delta i_C)^2P_{rr}
+\hat r_C^2R_{\Delta i}+\alphab^2R_q.
\label{eq:c_covariance}
\end{equation}
Here $R_{\Delta v}=2R_v$, $R_{\Delta i}=2R_i$, and $R_q$ follows from current-noise variance, the ADC interval, and the integration length. Equation~\eqref{eq:c_covariance} propagates uncertainty; a biased $(r_C-\hat r_C)\Delta i_C$ remains bounded model mismatch rather than white noise.

\begin{table}[t]
\caption{Decoupled Observation Rows}
\label{tab:observations}
\centering
\footnotesize
\begin{tabular}{p{0.9cm}p{2.35cm}p{2.1cm}p{1.35cm}}
\toprule
Row & Measurement & Regression & Dominant parameter \\
\midrule
Edge & $\hat v_T^- -\hat v_T^+$ & $k_R I_\Sigma r_C+\nu_R$ & $r_C$ \\
Charge & $\Delta v_T-\hat r_C\Delta i_C$ & $(q/C_b)\alphab+\nu_C$ & $\alphab$ \\
\bottomrule
\end{tabular}
\end{table}

\subsection{Validity and Update Policy}

An observation is accepted only in CCM with a valid timestamp, topology state, and complete safe window. Denoting these prechecks by $\mathcal S_C$ and $\mathcal S_R$, the common evaluation stream implements
\begin{equation}
\begin{aligned}
\mathcal V_C&=\mathrm{CCM}\wedge\mathcal S_C\wedge(n_C>0),\\
\mathcal V_R&=\mathrm{CCM}\wedge\mathcal S_R\wedge
(I_\Sigma\ge I_{\Sigma,\mathrm{gate}}),
\end{aligned}
\label{eq:gates}
\end{equation}
with the predeclared $I_{\Sigma,\mathrm{gate}}=0.12$~A. Charge acceptance uses the CCM/safe-window flag and a nonzero accepted-charge count; the implementation has no separate numeric $q_{\mathrm{gate}}$. An NIS threshold of 9 may reject additional \TSKF{} outliers. A failed gate holds the corresponding estimate. The evaluated-envelope minima in Section~\ref{sec:identifiability} are proof constants, not online thresholds.

Since the capacitance regressor is $h_C=q/C_b$, low-charge rows inherently contribute little information. The positive information guarantee is therefore restricted to the evaluated valid-CCM envelope rather than asserted for arbitrarily small $q$.

\begin{figure*}[t]
\centering
\includegraphics[width=0.97\textwidth]{fig02_observation_framework.pdf}
\caption{Core topology-synchronous decoupled observation framework. The edge-fit window $T_{w,R}$ and charge-safe window $T_{w,C}$ remain distinct; the downstream recursive estimators use standard scalar update principles.}
\label{fig:framework}
\end{figure*}
